\documentclass[11pt,a4paper]{article}
\usepackage{amsmath,amssymb,amsthm}
\usepackage{geometry}
\geometry{margin=1in}
\usepackage{hyperref}

\newtheorem{definition}{Definition}[section]
\newtheorem{theorem}[definition]{Theorem}
\newtheorem{proposition}[definition]{Proposition}
\newtheorem{remark}[definition]{Remark}

\title{\bfseries Harmonic Sine Transform as a Greedy Dirichlet Sub‑Sum Transform\\
\large Connecting Greedy Harmonic Decompositions to the Dirichlet Eta Function}
\author{}
\date{}

\begin{document}
\maketitle

\begin{abstract}
We reformulate the Harmonic Sine Transform (HST) of the harmonic–categorical framework as a \emph{Greedy Dirichlet Sub‑Sum Transform (GDST)}, directly linking the greedy harmonic expansion of angles to the Dirichlet eta function \(\eta(s)=(1-2^{1-s})\zeta(s)\).  Each target angle \(\theta\in[0,\pi]\) selects, via the greedy digits \(\delta_n(\theta)\in\{0,1\}\), a sub‑sum of the shifted Dirichlet eta series.  The resulting one‑parameter family of Dirichlet series \(E(\theta,s)\) and its signed variant \(E_\pm(\theta,s)\) provide an analytic continuation of the HST into the critical strip \(\operatorname{Re}s>0\).  This reformulation anchors the HST in classical analytic number theory and opens new pathways for studying the Riemann zeros through the lens of greedy approximations.
\end{abstract}

\section{Introduction}
The Harmonic Sine Transform, originally defined as an isometric isomorphism
\[
\mathcal{H}: L^2[0,\pi] \longrightarrow \ell^2(\mathcal{B}),\qquad
\mathcal{H}f(\psi)=\langle f,\psi\rangle_{L^2},
\]
maps square‑integrable functions on \([0,\pi]\) to sequences of wavelet coefficients obtained from the greedy harmonic tree.  The present note recasts this transform in the language of Dirichlet series.  The central object is the \emph{greedy Dirichlet sub‑sum}
\[
E(\theta,s) \;=\; \sum_{n=0}^{\infty} \frac{\delta_n(\theta)}{(n+2)^s},
\qquad \operatorname{Re}s>1,
\]
where \(\delta_n(\theta)\in\{0,1\}\) are the greedy digits determined by the harmonic fractions \(\pi/(n+2)\).  The signed version
\[
E_\pm(\theta,s) \;=\; \sum_{n=0}^{\infty} (-1)^n \frac{\delta_n(\theta)}{(n+2)^s},
\qquad \operatorname{Re}s>0,
\]
connects directly to the Dirichlet eta function.  We call the mapping
\[
\Theta: \theta \longmapsto E(\theta,\cdot)
\]
the \emph{Greedy Dirichlet Sub‑Sum Transform (GDST)}.

\section{Preliminaries: Greedy Harmonic Expansion}
We recall the essential facts from the companion paper \cite{Geere2026GHeta}.  The harmonic fractions are \(\alpha_n=\pi/(n+2)\) for \(n\ge 0\).  For \(\theta\in[0,\pi]\) the greedy algorithm produces binary digits
\[
\delta_n(\theta)=
\begin{cases}
1, & \text{if }\theta-\theta_n\ge\alpha_n,\\
0, & \text{otherwise},
\end{cases}
\qquad
\theta_{n+1}=\theta_n+\delta_n(\theta)\alpha_n,\quad \theta_0=0.
\]
The series expansion
\[
\frac{\theta}{\pi} = \sum_{n=0}^{\infty}\frac{\delta_n(\theta)}{n+2}
\]
converges absolutely, and the map \(\theta\mapsto\delta_n(\theta)\) is non‑decreasing with threshold angles \(\theta_n^*\) satisfying \(\delta_n(\theta)=\Theta(\theta-\theta_n^*)\).

\section{Greedy Dirichlet Sub‑Sum as a Transform}
Let \(s\in\mathbb{C}\) with \(\operatorname{Re}s>1\).  The greedy Dirichlet sub‑sum is the family of Dirichlet series
\[
E(\theta,s) = \sum_{n=0}^{\infty} \frac{\delta_n(\theta)}{(n+2)^s}
           = \sum_{\substack{n\ge 0\\\delta_n(\theta)=1}} \frac{1}{(n+2)^s}.
\]
For fixed \(s\) the map \(\theta\mapsto E(\theta,s)\) is a non‑decreasing step function with jumps of magnitude \((n+2)^{-s}\) at the threshold angles \(\theta_n^*\).  Thus
\[
E(\theta,s) = \sum_{n: \theta_n^*\le\theta} \frac{1}{(n+2)^s},
\]
which expresses \(E(\theta,s)\) as a cumulative distribution function of the harmonic sequence weighted by the zeta kernel.

\begin{definition}[Greedy Dirichlet Sub‑Sum Transform]\label{def:GDST}
For \(\operatorname{Re}s>1\) the \textbf{Greedy Dirichlet Sub‑Sum Transform} (GDST) of the angle \(\theta\) is the holomorphic function
\[
\Theta(\theta)(s) \;=\; E(\theta,s),\qquad s\in\{z:\operatorname{Re}z>1\}.
\]
\end{definition}

\section{The Signed Version and the Dirichlet Eta Function}
Much of the analytic content of the zeta function resides in the critical strip.  While \(E(\theta,s)\) converges only for \(\operatorname{Re}s>1\), its signed analogue
\[
E_\pm(\theta,s) = \sum_{n=0}^{\infty} (-1)^n \frac{\delta_n(\theta)}{(n+2)^s}
\]
converges conditionally for \(\operatorname{Re}s>0\) because the alternating signs together with the bounded cumulative counts \(\sum_{n\le N}(-1)^n\delta_n(\theta) = O(\ln N)\) provide the necessary cancellation \cite{Geere2026GHeta}.

\begin{proposition}[Connection to the Dirichlet eta function]\label{prop:eta}
For any \(\theta\in[0,\pi]\) the full alternating shifted series decomposes as
\[
\sum_{n=0}^{\infty}\frac{(-1)^n}{(n+2)^s}
   \;=\; E_\pm(\theta,s) \;+\; \Omega_\pm(\theta,s),
\]
where \(\Omega_\pm(\theta,s)=\sum_{n:\delta_n(\theta)=0}(-1)^n(n+2)^{-s}\) is the complement.
At \(\theta=\pi\) we obtain
\[
1-\eta(s) \;=\; E_\pm(\pi,s) \;+\; \sum_{n:\delta_n(\pi)=0}\frac{(-1)^n}{(n+2)^s},
\]
using the identity \(\sum_{n=0}^\infty (-1)^n (n+2)^{-s}=1-\eta(s)\).
Thus the Dirichlet eta function is partitioned by the greedy selection into a selected sub‑sum and its complement.
\end{proposition}

This partition is the core of the GDST: each angle \(\theta\) defines a sieve on the integers (shifted by two) that selects exactly those \(n\) for which \(\delta_n(\theta)=1\).  The signed selected sum \(E_\pm(\theta,s)\) becomes a one‑parameter family interpolating between \(0\) (at \(\theta=0\)) and a value close to \(1-\eta(s)\) (at \(\theta=\pi\)).

\section{Harmonic Sine Transform via the GDST}
In the original HST, the wavelet coefficients \(\mathcal{H}f(\psi_\varepsilon)\) decompose a function \(f\in L^2[0,\pi]\) into sums over the greedy cylinder intervals.  Each Haar wavelet \(\psi_\varepsilon\) is supported on an interval defined by a finite prefix of the digits \(\delta_n\).  Therefore the HST can be re‑expressed in terms of the digit functions \(\delta_n\) themselves.

The Dirichlet series associated with a wavelet \(\psi_\varepsilon\) of depth \(m\) is
\[
D_{\psi_\varepsilon}(s) = \sum_{\psi\in\mathcal{B}}\frac{\langle\psi_\varepsilon,\psi\rangle}{(|\psi|+2)^s}
\]
but a more direct path uses the identity
\[
\delta_n(\theta) = \Theta(\theta-\theta_n^*),
\]
so that
\[
E(\theta,s) = \sum_{n=0}^{\infty} \frac{\Theta(\theta-\theta_n^*)}{(n+2)^s}.
\]
For a smooth test function \(f\in L^2[0,\pi]\), define its \textbf{GDST coefficients} by
\[
\widehat{f}(s) \;=\; \int_0^\pi f(\theta)\, E(\theta,s)\, d\theta,\qquad \operatorname{Re}s>1.
\]
This is precisely the form of the HST Dirichlet series \(D_f(s)\) introduced in the earlier programme, but now with \(E(\theta,s)\) directly linked to the greedy selection of the Dirichlet eta series.

\section{Asymptotic Approximation of \(\theta\) by Greedy Eta Sums}
The companion paper \cite{Geere2026GHeta} shows that the cumulative count
\[
D(N,\theta) = \sum_{n=0}^{N} \delta_n(\theta)
\]
satisfies
\[
D(N,\theta) = \frac{\theta}{\pi}\,\ln(N+2) + O(1).
\]
Consequently, for \(\operatorname{Re}s>1\),
\[
E(\theta,s) = \frac{\theta}{\pi}\bigl(\zeta(s)-1\bigr) \;+\; \text{absolutely convergent correction}.
\]
Thus the scaled family \(\frac{\pi}{\theta}E(\theta,s)\) provides a regularised approximation to the shifted zeta function \(\zeta(s)-1\), with the angle \(\theta\) acting as a parameter that tunes the density of selected indices.

In the limit \(\theta\to\pi\) the greedy algorithm selects most indices, and \(E(\pi,s)\) approximates \(\zeta(s)-1\) up to a sparse set of omitted terms (the indices where \(\delta_n(\pi)=0\)).  Conversely, at \(\theta=\pi/2\), the greedy selection picks only the very first term \(n=0\) because \(\alpha_0=\pi/2\) already reaches the half angle; all higher terms are excluded.  This behaviour allows the GDST to probe the Dirichlet eta series at different “scales” determined by the target angle.

\section{Discussion and Relation to the Riemann Hypothesis}
The reformulation of the HST as a Greedy Dirichlet Sub‑Sum Transform places the harmonic–categorical approach squarely within the analytic theory of the Riemann zeta function.  The determinant identity for the transfer operator \(\mathcal{L}_s\) is mirrored by the family of Dirichlet series \(E(\theta,s)\), whose behaviour in the critical strip is governed by the eta function.  The failure of the finite‑dimensional spectral correspondence (Programme~VII) suggests that a direct limit of truncated matrices is not the correct route to a proof of the Riemann Hypothesis.  However, the GDST viewpoint opens alternative avenues:

\begin{enumerate}
\item \textbf{Spectral asymptotics.}  The correlation kernel \(K(n,m)\) of the digit functions controls the \(L^2\) fluctuations of \(E(\theta,s)\) and may encode the distribution of zeta zeros through its spectral decomposition.
\item \textbf{Characteristic angles of zeros.}  Proposition~7.3 of \cite{Geere2026GHeta} shows that every non‑trivial zero \(\rho\) of \(\zeta(s)\) defines a \emph{balance angle} \(\theta_{\text{bal}}(\rho)\in[0,\pi]\) satisfying \(\operatorname{Re}E_\pm(\theta_{\text{bal}},\rho)=1/2\).  A detailed study of these characteristic angles could unveil a new symmetry of the zeta zeros.
\item \textbf{Integral representations.}  The generating function \(G(\theta,t)=\sum_n\delta_n(\theta)e^{-(n+2)t}\) admits a representation via the thresholds \(\theta_n^*\) and connects to the Mellin transform of the eta function.
\end{enumerate}

None of these observations constitutes a proof of the Riemann Hypothesis; the core obstacle—constructing a self‑adjoint operator whose spectrum matches the zeta zeros—remains.  But the GDST formulation provides a \(C^*\)-algebraic framework (functions on the greedy tree) that may be amenable to Connes’ spectral‑triple techniques.

\section{Conclusion}
We have defined the Harmonic Sine Transform as a Greedy Dirichlet Sub‑Sum Transform, linking the geometric greedy decomposition of angles to the analytic theory of the Dirichlet eta function.  The mapping \(\theta\mapsto E(\theta,s)\) is a rigorously defined one‑parameter family of Dirichlet series that interpolates within the shifted zeta function and, via its signed extension \(E_\pm(\theta,s)\), extends into the critical strip.  This perspective unifies the earlier harmonic–categorical programmes with classical analytic number theory and suggests new strategies for approaching the Riemann Hypothesis.

\begin{thebibliography}{9}
\bibitem{Geere2026GHeta}
V.~Geere, \emph{The Greedy Harmonic Decomposition and the Dirichlet Eta Function}, March 2026.
\url{https://victorgeere.co.za/maths/greedy-harmonic-eta.html}
\bibitem{Geere2026a}
V.~Geere, \emph{On the Correlation Structure of a Greedy Harmonic Decomposition of the Sine Function}, March 2026.
\bibitem{Geere2026b}
V.~Geere, \emph{A Harmonic Reconstruction of the Sine Function and Its Relation to the Riemann Zeta Function}, 2026.
\bibitem{HSTprogramme}
V.~Geere et al., \emph{A Harmonic Sine Transform and the Riemann Zeta Function}, 2026.
\end{thebibliography}

\end{document}